\begin{proposition}[Semistable limits of the bielliptic surfaces]\hypertarget{tag-01R6}{}\label{fam:src:E.2.1}\label{fam:st:E.2.1}
\StatementRecord{fam:src:E.2.1}{proved}{01R6}{the definitions and constructions stated in the proof}

After the auxiliary base change $w=s^6$, each of the two families of
\emph{reduced} special surfaces $S_j(w)$ has a proper flat
semistable filling by explicit finite free quotients. The central
surface for $j=1$ is a cycle of four ruled surfaces over an elliptic
curve; that for $j=2$ is a cycle of three. They have
\[
\begin{array}{c|ccccc|c}
 &H_0&H_1&H_2&H_3&H_4&\chi\\\hline
 j=1&\Z&\Z\oplus\Z/3&\Z^5&\Z^9&\Z^4&0\\
 j=2&\Z&\Z\oplus\Z/2&\Z^4&\Z^7&\Z^3&0
\end{array}
\]
with integral coefficients. Their fundamental groups are
$\Z^2\rtimes_{R_3}\Z$ and $\Z^2\rtimes_{R_4}\Z$, respectively,
where
\[
 R_3=\begin{pmatrix}0&-1\\1&-1\end{pmatrix},\qquad
 R_4=\begin{pmatrix}0&-1\\1&0\end{pmatrix}.
\]
\end{proposition}



\begin{proof}
\begin{enumerate}[label=\arabic*.,leftmargin=*]
\item The following product presentations specify the comparison
with the marked central surfaces. Let
$\rho=e^{\pi i/3}$,
\[
 E_j=\C/(\Z+\kappa_j\Z),\quad
 F_1=\C/(\Z(\rho-2)+\Z(-3)),\quad
 F_2=\C/(\Z(i-1)+\Z(-2)),
\]
where $\kappa_1=\beta_c(z_1)+2\mu(z_1)$ and
$\kappa_2=\beta_c(z_2)+3\mu(z_2)$ have negative imaginary parts.
The holomorphic maps to the central tori are induced by
$L_1(\xi,\eta)=(\eta,\xi-\eta/3)$ and
$L_2(\xi,\eta)=(\eta,\xi-\eta/2)$. Their kernels are generated by
\[
 k_1=(2/3,\rho),\qquad k_2=(1/2,i).
\]
The residual affine generators lift to
\begin{equation}
 h_1(\xi,\eta)=(\xi+\kappa_1/3,(\rho-1)\eta),\qquad
 h_2(\xi,\eta)=(\xi-\kappa_2/4,-i\eta).
\label{pub:I:eq:E8}
\end{equation}
These formulas are verified by substituting the four product periods
in $L_j$; the lattice indices are three and two. Moreover
$(\rho-1)\rho-\rho$ and $(-i)i-i$ are periods of the displayed
$F_j$, so the generators commute. They give free groups
$G_1=C_3\times C_3$ and $G_2=C_2\times C_4$. Indeed, equality of
the first coordinates of $k_j^a h_j^b(\xi,\eta)$ and $(\xi,\eta)$
forces $b=0$ from the coefficient of the nonreal period $\kappa_j$,
and then $a=0$ from the real period. Thus the composite quotient
maps are exactly the degree-nine and degree-eight covers of $S_j$,
with the stated signs.

\item Use the multiplicative coordinate $v=e^{-2\pi i\xi}$ on $E_j$.
Its Tate parameter is
\[
       q_j=e^{-2\pi i\kappa_j}=A_jw^2=A_js^{12},
       \qquad A_j=e^{-2\pi i\kappa_j^{\mathrm{part}}}\ne0.
\]
Here $A_j$ denotes the displayed scalar constant within this proof. Choose specifically
$A_j^{1/12}=\exp(-2\pi i\kappa_j^{\mathrm{part}}/12)$ and put
$r_j=A_j^{1/12}s$. This choice preserves the fractional period
translations in \eqref{pub:I:eq:E8} exactly. The regular Tate model for $q_j=r_j^{12}$
has a central twelve-gon. We recall its charts to make the extension
explicit. An infinite chain has node charts
$x_i y_i=r_j$; on a branch overlap one uses
$x_{i+1}=y_i^{-1}$ and $y_{i+1}=r_jy_i$. The coordinate
$r_j^i x_i$ on a general fibre is unchanged on overlaps.
The shift $i\mapsto i+12$ acts on that coordinate by $r_j^{12}$.
Quotienting a tubular neighbourhood of the chain by this shift gives
the proper twelve-gon family, and on $r_j\ne0$ it is
$\C^*/q_j^\Z$. Compactness follows from twelve projective components
and a compact annular fundamental domain on a general fibre;
the node charts give flatness and smoothness of its total space.

\item Multiplication by $q_1^{1/3}=r_1^4$ shifts the component index by four.
Multiplication by $q_2^{-1/4}=r_2^{-3}$ shifts it by minus three.
Hence the maps $h_j$ in \eqref{pub:I:eq:E8} extend, together with their rotations
on $F_j$. The maps $k_j$ multiply $v$ by roots of unity and translate
$F_j$, and extend to every projective component. Their relations
hold on the dense set $s\ne0$ and consequently on all charts.
The action is free on the central fibre as well. A nontrivial power
of $h_j$ moves every component and every node. A nontrivial power
of $k_j$ translates $F_j$ by a nonzero torsion point, including
above the nodes where the multiplicative action alone fixes a point.
The finite free quotient is therefore a smooth total space, proper
over $s$, with a reduced normal-crossing central surface. Flatness
descends along the finite etale quotient. The component orbits have
cardinality three and four, leaving four and three components.
Quotienting an individual $\mathbb P^1\times F_j$ by $k_j$ is a
ruled surface over $F_j/\langle\rho\rangle$ or
$F_j/\langle i\rangle$, respectively. Its two distinguished
sections have the same integral homology class: the corresponding
line bundle has finite order and degree zero on the elliptic base.

\item For the topology, unwrap the component cycle into an infinite chain
and then pass to the universal cover $\C$ of its elliptic base.
The resulting union of simply connected ruled pieces, with
simply connected intersections and a tree as incidence graph, is
simply connected by van Kampen. Its deck group is generated by
the enlarged lattice of translations on $F_j$ and one step round
the component cycle. The latter conjugates the lattice by the
order-three or order-four rotation. This gives the asserted
semidirect products; inverses of the displayed rotation matrices
give isomorphic groups by reversing the cycle generator.

\item For homology, the normalization has $n=4$ or $3$ components, each
with ranks $(1,2,2,2,1)$, and $n$ elliptic double curves, each with
ranks $(1,2,1)$. There are no triple points. The mapping-cylinder
sequence for the normalization and the double curves applies
integrally. In degree one the incidence matrix is a twisted cycle:
its cokernel is $\operatorname{coker}(R_j-I)$ and its kernel is
zero, because $R_j$ has no eigenvalue one. The degree-zero incidence
matrix contributes an additional $\Z$ to $H_1$. In degree two the
two section classes on a component agree, so the incidence matrix
is the ordinary cycle matrix on their classes. It has kernel
$\Z$ and primitive image of rank $n-1$. Thus $H_2=\Z^{n+1}$,
$H_3=\Z^{2n+1}$, and $H_4=\Z^n$. Finally
$\operatorname{coker}(R_3-I)=\Z/3$ and
$\operatorname{coker}(R_4-I)=\Z/2$. 
\end{enumerate}
\end{proof}